\chapter{Buffervaten V-101 en V-102}
\label{h:buffers}
Een buffervat slaat gas op door de druk te laten variëren tussen $P_{\max}$ en $P_{\min}$. De hoeveelheid gas die
je eruit kunt halen, volgt uit het verschil in dichtheid tussen beide drukken. Alles geldt \textbf{per unit}.

\section{Opgave}
\begin{center}
\begin{tabular}{p{3.3cm}p{5.6cm}p{5.6cm}}
\toprule
 & V-101 (lagedrukbuffer) & V-102 (hogedrukbuffer)\\
\midrule
functie & ITM-productie opvangen als K-101 stilstaat of anders draait & drukstoten naar de backbone dempen\\
druk & $\VAR{n(I.P_V101_max,0)} \rightarrow \VAR{n(I.P_V101_min,0)}$\,barg & $\VAR{n(I.P_K101_uit,0)} \rightarrow \VAR{n(I.P_backbone,1)}$\,barg\\
$P_{\max}$, $P_{\min}$ [kPa] & \VAR{n(R.P_V101_max if A else R.V101.pmax,1)}; \VAR{n(R.P_V101_min if A else R.V101.pmin,1)} & \VAR{n(R.P_V102_max if A else R.V102.pmax,1)}; \VAR{n(R.P_V102_min if A else R.V102.pmin,1)}\\
buffertijd & \VAR{n(I.t_V101,1)}\,h & \VAR{n(I.t_V102,0)}\,s\\
te leveren massa $\Delta m$ & $\VAR{n(I.m_H2,0)}\cdot\VAR{n(I.t_V101,1)} = \VAR{n(I.m_H2*I.t_V101,1)}$\,kg & $\VAR{n(I.m_H2,0)}\cdot\VAR{n(I.t_V102,0)}/3600 = \VAR{n(I.m_H2*I.t_V102/3600,3)}$\,kg\\
temperatuur & 25\degC & 25\degC\\
ontwerpfactor $n$ & \VAR{n(I.n_V101,1)} (polytroop) & $k = \VAR{n(R.k_H2 if A else C.k_H2,3)}$ (adiabatisch)\\
\bottomrule
\end{tabular}
\end{center}
\begin{waarom}[de buffertijden]
V-101: een half uur geeft operators tijd om bij een storing van K-101 in te grijpen, voordat de ITM moet afschakelen
of moet afblazen. V-102: 30\,s is genoeg om de regeling van PCV-101 en de opstart of stop van een compressor op te
vangen zonder drukdip in de backbone. Beide zijn ontwerpaannames van het model.
\end{waarom}

\section{Methode}
\begin{methode}[– Cengel H3 (toestandsvergelijking) en H4/H7 (polytroop proces)]
Massa in het vat: $m = \rho V = \dfrac{P V}{Z R T}$. Levert het vat $\Delta m$ bij \emph{constante} temperatuur
(isotherm), dan geldt
\[
\Delta m = V\,(\rho_{\max}-\rho_{\min}) \quad\Rightarrow\quad V_\mathrm{iso} = \frac{\Delta m}{\rho_{\max}-\rho_{\min}}.
\]
Bij snel leeg laten lopen koelt het gas in het vat af. Het gas dat achterblijft expandeert dan polytroop
($P v^n = $ constant, $1 \le n \le k$), heeft een hogere dichtheid en levert dus \emph{minder} massa. Het vat moet
groter zijn: $V = n\cdot V_\mathrm{iso}$.
\end{methode}
\kop{Afleiding van de factor $n$}
Met $T/T_1 = (P/P_1)^{(n-1)/n}$ (polytroop ideaal gas, Cengel H7) is de massa in het vat
$m(P) = \frac{V P_1}{R T_1}\left(\frac{P}{P_1}\right)^{1/n}$. Dus
\[
\Delta m_\mathrm{poly} = \frac{V P_1}{R T_1}\left[1 - \left(\frac{P_2}{P_1}\right)^{1/n}\right],\qquad
\Delta m_\mathrm{iso} = \frac{V P_1}{R T_1}\left[1 - \frac{P_2}{P_1}\right].
\]
Voor een kleine drukzwaai is $1-(P_2/P_1)^{1/n} \approx \frac{1}{n}(1-P_2/P_1)$, dus $V_\mathrm{poly} \approx n\,V_\mathrm{iso}$.
De exacte verhouding is voor V-101 $\frac{1-0{,}700}{1-0{,}700^{1/1{,}2}} = 1{,}17$ en voor V-102
$\frac{1-0{,}947}{1-0{,}947^{1/1{,}405}} = 1{,}39$. De vuistregel ``$\times n$'' is dus iets conservatief (V-101)
of vrijwel exact (V-102).
\begin{waarom}[$n = 1{,}2$ voor V-101 en $n = k$ voor V-102]
V-101 loopt in een half uur leeg. Er is tijd voor warmte-uitwisseling met de wand, dus het proces ligt tussen
isotherm ($n=1$) en adiabatisch ($n=k$); 1,2 is een gangbare keuze. V-102 loopt in 30\,s leeg. Dat is te snel
voor warmte-uitwisseling, dus het proces is adiabatisch: $n = k = 1{,}405$.
\end{waarom}

%% if A
\section{V-101 met CoolProp}
\begin{coolprop}
$\rho_{\max} = $ \texttt{PropsSI("D","T",298.15,"P",\VAR{n(R.P_V101_max*1000,0)},"Hydrogen")} $= \VAR{n(R.rho_V101_max,4)}$\,kg/m$^3$
($Z = \VAR{n(R.Z_V101_max,4)}$);\\
$\rho_{\min} = $ \texttt{PropsSI("D","T",298.15,"P",\VAR{n(R.P_V101_min*1000,0)},"Hydrogen")} $= \VAR{n(R.rho_V101_min,4)}$\,kg/m$^3$
($Z = \VAR{n(R.Z_V101_min,4)}$).\\
Cengel-check (ideaal gas, $\rho = P/(RT)$): $\VAR{n(R.P_V101_max,1)}/(4{,}124\cdot298{,}15) = \VAR{n(R.P_V101_max/(4.124*298.15),4)}$ en
$\VAR{n(R.P_V101_min/(4.124*298.15),4)}$\,kg/m$^3$. Ideaal gas overschat de dichtheid met ongeveer $Z-1$.
\end{coolprop}
\[
V_\mathrm{iso} = \frac{\VAR{n(R.m_buf_V101,1)}}{\VAR{n(R.rho_V101_max,4)} - \VAR{n(R.rho_V101_min,4)}} = \VAR{n(R.V_V101_iso,2)}\,\mathrm{m^3},\qquad
V_{\text{V-101}} = \VAR{n(I.n_V101,1)}\cdot\VAR{n(R.V_V101_iso,2)} = \mathbf{\VAR{n(R.V_V101,1)}\,m^3}.
\]
\section{V-102 met CoolProp}
\[
\rho_{\max} = \VAR{n(R.rho_V102_max,4)},\quad \rho_{\min} = \VAR{n(R.rho_V102_min,4)}\,\mathrm{kg/m^3}\ (Z = \VAR{n(R.Z_V102_max,3)}),\quad
V_\mathrm{iso} = \frac{\VAR{n(R.m_buf_V102,3)}}{\VAR{n(R.rho_V102_max - R.rho_V102_min,4)}} = \VAR{n(R.V_V102_iso,3)}\,\mathrm{m^3},
\]
\[
V_{\text{V-102}} = k\cdot V_\mathrm{iso} = \VAR{n(R.k_H2,4)}\cdot\VAR{n(R.V_V102_iso,3)} = \mathbf{\VAR{n(R.V_V102,2)}\,m^3}.
\]
Hier is $k = c_p/c_v$ bij 25\degC\ en 1\,atm: \texttt{PropsSI("CPMASS",…)/PropsSI("CVMASS",…)} $= \VAR{n(R.k_H2,4)}$
(Cengel A-2: 1,405).
%% else
\section{V-101 met ideaal gas en viriaalcorrectie}
\begin{tabelkader}[A-1 + viriaalcoëfficiënt]
$R_{\HH} = 4{,}124\kJkgK$; tweede viriaalcoëfficiënt van \HH\ bij 25\degC: $B = 14{,}8$\,cm$^3$/mol $= 14{,}8\cdot10^{-6}$\,m$^3$/mol.
\end{tabelkader}
\begin{align*}
RT &= 4{,}124\cdot 298{,}15 = \VAR{n(R.RT,2)}\,\mathrm{kJ/kg}, \qquad R_uT = 8{,}314\cdot298{,}15 = \VAR{n(R.RuT,1)}\,\mathrm{J/mol},\\
Z_{\max} &= 1 + \frac{B\,P_{\max}}{R_u T} = 1 + \frac{14{,}8\cdot10^{-6}\cdot\VAR{n(R.V101.pmax*1000,0)}}{\VAR{n(R.RuT,1)}} = \VAR{n(R.V101.zmax,4)},\qquad
Z_{\min} = \VAR{n(R.V101.zmin,4)},\\
\rho_{\max} &= \frac{P_{\max}}{Z_{\max} R T} = \frac{\VAR{n(R.V101.pmax,1)}}{\VAR{n(R.V101.zmax,4)}\cdot\VAR{n(R.RT,2)}} = \VAR{n(R.V101.rho_max,4)}\,\mathrm{kg/m^3},\qquad
\rho_{\min} = \VAR{n(R.V101.rho_min,4)}\,\mathrm{kg/m^3},\\
V_\mathrm{iso} &= \frac{\VAR{n(R.V101.m,1)}}{\VAR{n(R.V101.rho_max,4)} - \VAR{n(R.V101.rho_min,4)}} = \VAR{n(R.V101.isoZ,2)}\,\mathrm{m^3}
\quad(\text{zonder }Z\text{: }\VAR{n(R.V101.iso,2)}\,\mathrm{m^3}),\\
V_{\text{V-101}} &= \VAR{n(I.n_V101,1)}\cdot\VAR{n(R.V101.isoZ,2)} = \mathbf{\VAR{n(R.V101.V,1)}\,m^3}.
\end{align*}
\section{V-102}
\begin{align*}
Z_{\max} &= \VAR{n(R.V102.zmax,4)},\quad Z_{\min} = \VAR{n(R.V102.zmin,4)},\quad
\rho_{\max} = \VAR{n(R.V102.rho_max,4)},\quad \rho_{\min} = \VAR{n(R.V102.rho_min,4)}\,\mathrm{kg/m^3},\\
V_\mathrm{iso} &= \frac{\VAR{n(R.V102.m,3)}}{\VAR{n(R.V102.rho_max - R.V102.rho_min,4)}} = \VAR{n(R.V102.isoZ,3)}\,\mathrm{m^3},\qquad
V_{\text{V-102}} = 1{,}405\cdot\VAR{n(R.V102.isoZ,3)} = \mathbf{\VAR{n(R.V102.V,2)}\,m^3}.
\end{align*}
%% endif
\begin{resultaat}
V-101 $\approx \VAR{n(R.V_V101,0)}$\,m$^3$ per unit en V-102 $\approx \VAR{n(R.V_V102,1)}$\,m$^3$ per unit. V-101 is groot
omdat een half uur productie bij een drukzwaai van maar 9\,bar moet passen. Een grotere drukzwaai of een kortere
buffertijd maakt het vat kleiner, maar de ondergrens van 20\,barg is de minimale zuigdruk van K-101 en de
bovengrens van 29\,barg ligt net onder de leverdruk van de ITM (30\,barg).
\end{resultaat}
\begin{waarom}[$Z$ meenemen]
Bij 30 tot 70\,bar wijkt \HH\ 2 tot 4\,\% af van ideaal gas. In een dichtheid\emph{verschil} kan die afwijking groter
uitvallen, omdat $Z$ bij beide drukken anders is. Zonder $Z$ zou V-101 met
\VAR{n((1 - RB.V101.iso / (R.V_V101_iso if A else R.V101.isoZ))*100,1)}\,\% te klein
uitvallen.
\end{waarom}
